\section{Introduzione Generale}
L'esercizio assegnatomi per il completamento del corso di Laboratorio di Algortimi e Strutture Dati viene riportato come segue:
\begin{itemize}
	\item {\textit  {"Vogliamo confrontare vari modi per costruire alberi di ricerca: normali ABR, alberi Rosso-Neri, alberi AVL."}}
\end{itemize}
	Nel paragrafo sottostante verranno elencati i punti salienti della intera relazione.

\subsection{Cosa affronterò}

Si può suddividere la relazione in 4 macro sezioni:

\begin{itemize}
	\item \textbf {Teoria e Problema}: in questa sezione affronterò l'esercizio da risolvere sfruttando la teoria appresa durante il corso di Algoritmi e Strutture Dati. Si descriverà brevemente le strutture dati che verranno utilizzate, gli algoritmi usati su tali strutture ed i risultati attesi (secondo la teoria) che verranno poi messi a confronto con quelli sperimentali.
	\item \textbf{Documentazione del codice}: in questa sezione descriverò come il codice python è stato implementato.
	\item \textbf{Descrizione degli esperimenti}: in questa sezione descriverò i risultati sperimentali ottenuti con grafici e tabelle.
	\item \textbf{Analisi dei risultati sperimentali}: in questa sezione esporrò una analisi critica di tutti i risultati ottenuti effettuando così le mie conclusioni.
\end{itemize}

\subsection{Calcolatore e Ambienti}
Di seguito vengono riportate le specifiche del Calcolatore utilizzato per il conseguimento dei test sperimentali; ambienti e linguaggio usati per la stesura del programma e della relazione stessa.

\begin{itemize}
	\item \textbf{Specs del Calcolatore}:
		\begin{itemize}
			\item \textbf{CPU}: Chip Apple M1, 8-core
			\item \textbf{RAM}: 8 GB
			\item \textbf{SSD}: 512 GB
		\end{itemize}
	\item \textbf{Linguaggio di Programmazione}: Python
	\item \textbf{Ambienti}: VS-Code, TexStudio
\end{itemize}